%!TEX root = ../graduate-work.tex
\phantomsection
\section*{Abstract}
\addcontentsline{toc}{section}{Abstract}

Formal verification of neural networks---the mathematical proof that a network
satisfies given safety properties for all inputs within a specified domain---faces
a fundamental scalability bottleneck: bound-propagation algorithms (IBP, CROWN,
$\alpha$-CROWN) require weight matrices to be stored in their entirety on a single
GPU, which precludes verification of models with billions of parameters.

This thesis investigates the application of parallelism techniques developed for
training large language models to the problem of formal neural network verification.
Two approaches are implemented and experimentally evaluated within the
\texttt{auto\_LiRPA} / $\alpha,\beta$-CROWN framework:

\textbf{Tensor Parallelism (TP)}---sharding of weight matrices and linear-relaxation
coefficient matrices ($A$-matrices) across multiple GPUs. A peak-memory reduction of
$\approx\!2\times$ is achieved at $P=2$ GPUs; soundness of the computed bounds is
confirmed on VNN-COMP 2022 benchmarks (MNIST-FC).

\textbf{Fully Sharded Data Parallelism (FSDP)}---layer-wise sharding of weight
matrices only, with an \texttt{AllGather} collective before each layer.
The resulting bounds are bitwise identical to the single-GPU baseline: baseline
memory (weight storage) is reduced by exactly 50\% (factor $1/P$ at $P{=}2$),
and peak memory during \texttt{compute\_bounds} by 17--39\% on MLPs. FSDP is integrated with
complete verification ($\beta$-CROWN + Branch-and-Bound) and extended to
convolutional layers (\texttt{BoundConv}); correctness is verified on
CIFAR-100 ResNet medium/large (VNN-COMP 2024), with ResNet-large yielding a
complete verification result (unsat).

Experiments establish that in $\alpha$-CROWN+BaB mode the memory bottleneck
is the per-neuron alpha tensors rather than weight matrices. The applicability
boundaries of each approach are characterised and directions for future work are formulated.
